\documentclass[10pt,a4paper]{article}
\usepackage[margin=2.2cm]{geometry}
\usepackage[T1]{fontenc}
\usepackage{lmodern}
\usepackage{microtype}
\usepackage{amsmath,amssymb}
\usepackage{xcolor}
\usepackage{booktabs}
\usepackage{longtable}
\usepackage{enumitem}
\usepackage[font=small,labelfont=bf]{caption}
\usepackage[colorlinks=true,linkcolor=blue!50!black,urlcolor=blue!50!black]{hyperref}

\setlength{\parindent}{0pt}
\setlength{\parskip}{0.4em}
\setlist{nosep,leftmargin=*,topsep=3pt}
\renewcommand{\arraystretch}{1.12}

\newcommand{\nbar}{\bar n}
\newcommand{\dtil}{\tilde\delta}
\newcommand{\B}[1]{\textbf{#1}}

\newcommand{\bestclassical}{0.719}
\newcommand{\bestclassicalname}{$L{+}F{+}G$}
\newcommand{\bestph}{0.666}
\newcommand{\lonly}{0.685}
\newcommand{\phbelow}{0.358}
\newcommand{\lonlyxbelow}{0.474}
\newcommand{\bestxbelow}{0.473}
\newcommand{\phweak}{0.643}
\newcommand{\lonlyxweak}{0.659}
\newcommand{\bestxweak}{0.684}
\newcommand{\phmoderate}{0.784}
\newcommand{\lonlyxmoderate}{0.766}
\newcommand{\bestxmoderate}{0.819}
\newcommand{\phstrong}{0.878}
\newcommand{\lonlyxstrong}{0.841}
\newcommand{\bestxstrong}{0.899}


\title{\textbf{Point-process TDA: status and results}}
\date{\today}

\begin{document}
\maketitle

% ============================================================================
\section{Where things stand}

\begin{itemize}
  \item Results are now reported on a \textbf{single stratified test set}: 50{,}000 patterns, balanced over the five models and over four bands of ``distance from Poisson''. It replaces the three separately-named test sets used previously, which made the headline number depend on which set a reader happened to look at.
  \item The test set was assembled \emph{by selection} from patterns that already existed. Nothing was re-simulated, no features were recomputed and no model was retrained --- every number below comes from the predictions already on disk. See \texttt{docs/test\_set.md}.
  \item \textbf{Done:} the classical baselines (a neural network reading the standard summary functions), for parameter estimation and 5-way model classification, covering all 15 non-empty subsets of $\{L,F,G,J\}$; and the first persistence-homology (PH) features: persistence images of the DTM filtration ($k=5$), $H_0$ and $H_1$ separately and both together. Every number is an average over 10 training seeds.
  \item \textbf{Main result:} a classical model wins on every task, but the margin is concentrated almost entirely in one band. For classification the best classical model (\bestclassicalname) reaches \bestclassical\ against \bestph\ for the best PH model ($H_0{+}H_1$). Split by distance from Poisson, PH is \emph{ahead of the published $L$-only baseline} in the two most structured bands (\phmoderate\ vs.\ \lonlyxmoderate, and \phstrong\ vs.\ \lonlyxstrong) and loses the comparison entirely in the band where the classical test cannot detect a departure at all (\phbelow\ vs.\ \lonlyxbelow).
  \item \textbf{Important caveat on the PH numbers:} every PH run so far uses a \emph{single} filtration scale ($k=5$) and omits the point count $n(x)$ for classification. The configuration that won in the earlier report --- multi-scale $k=5,10,15$ --- has not yet been run. The PH column is therefore an ablation, not the method.
  \item \textbf{Not yet run:} multi-scale DTM ($k=10,15$); $n(x)$ included for classification; Betti curves and persistence landscapes; the Rips/alpha filtration; and the non-learned classical references (minimum-contrast fitting, envelope tests).
\end{itemize}

% ============================================================================
\section{Data}

\textbf{Models.} Five stationary models on the unit square: Poisson (complete spatial randomness, CSR), Thomas (clustered), nested Thomas (clusters of clusters), Mat\'ern type~II (hard-core, i.e.\ repulsive) and the log-Gaussian Cox process (LGCP). The expected number of points $\nbar$ ranges from 100 to 800.

\textbf{How the distance from Poisson is measured ($\dtil$).}
\begin{itemize}
  \item The classical way to test a pattern against Poisson is to compare its $L$-function ($L(r)-r$ is zero for Poisson) with its Poisson expectation, taking the largest standardised deviation over a range of radii (a global envelope test).
  \item For every parameter setting we compute that largest deviation for the model's \emph{theoretical} $L$-function, and divide it by the test's 5\% critical value at the same $\nbar$. That ratio is $\dtil$.
  \item So $\dtil = 0$ is exactly Poisson, and $\dtil \le 1$ means precisely: \emph{the classical 5\% test does not reject Poisson}. That is not a round number chosen after the fact --- it is the test's own critical value, fixed before any of these results were produced, and it is what the test set is cut on.
  \item At a fixed shape, $\dtil$ grows roughly like $\sqrt{\nbar}$: detecting a given departure takes points.
\end{itemize}

\textbf{The test set.} One table of 50{,}000 patterns, \textbf{2{,}500 in every cell} of (5 models $\times$ 4 bands):

\begin{center}\small
\begin{tabular}{@{}lll@{}}
\toprule
Band & $\dtil$ & Meaning\\
\midrule
\texttt{below}    & $<1$   & the classical test does not reject Poisson\\
\texttt{weak}     & 1--2   & just past the detection threshold\\
\texttt{moderate} & 2--4   & comfortably detectable\\
\texttt{strong}   & $\ge 4$ & the departure is unambiguous\\
\bottomrule
\end{tabular}
\end{center}

Equal cells mean the overall number is an \emph{explicit, stated weighting} of the four regimes rather than an accident of the prior the patterns were drawn from --- the previous headline set put roughly half its mass in \texttt{below}, so it was measuring the floor of the problem as much as the methods.

Two details that matter when reading the tables:
\begin{itemize}
  \item \textbf{Poisson is a class, not a band.} CSR patterns have $\dtil = 0$, so a pure $\dtil$ cut would put all of them in \texttt{below} and leave the other bands as 4-way problems with a different chance level. Instead each band gets the same number of CSR patterns, drawn without overlap, so chance stays at 0.20 in every band and the columns are comparable. A CSR pattern's band names the band it is \emph{paired with}.
  \item \textbf{$\nbar$ is recorded, not balanced.} Reaching a given $\dtil$ requires points, so $\dtil$ and $\nbar$ are entangled by construction and no selection can separate them --- Mat\'ern~II in particular cannot exceed $\dtil \approx 8$ at all, because its hard core is bounded by a packing constraint. Each cell is filled as evenly across the three $\nbar$ octaves as the available patterns allow and the realised mix is recorded per row.
\end{itemize}

\textbf{Training} is unchanged: 10{,}000 patterns per model from the prior, with a fixed 85/15 train/validation split assigned at generation and identical for every method and seed. The test set is disjoint from it.

\textbf{Error bars.} Tables report mean $\pm$ standard deviation over the 10 seeds. The test set pools patterns that share a parameter vector (some sources contribute hundreds of replicates at one setting), so the per-seed standard errors quoted in the machine-readable output are \emph{clustered} on the parameter vector: the 50{,}000 patterns resolve to 14{,}041 distinct settings, and treating them as 50{,}000 independent draws would understate the uncertainty.

% ============================================================================
\section{Methods compared}

\begin{itemize}
  \item \textbf{Classical baselines:} the 1-D CNN of Vihrs (2022), reading summary functions sampled on a radius grid, plus the point count $n(x)$. All 15 non-empty subsets of $\{L,F,G,J\}$ were tried, from a single function ($L$ alone is the published baseline) up to the full union $L{+}F{+}G{+}J$. $L$, $L{+}F{+}G{+}J$, $G$ and $F$ alone were additionally tried on two radius grids for $F/G/J$ ($r_{\max}=0.25$ and $0.08$); the tables report whichever grid won.
  \item \textbf{PH:} a CNN on the persistence image of the DTM filtration ($k=5$): $H_0$ or $H_1$ alone (separate models), or both as the two channels of one image. Image pixels are standardised per channel with training-set statistics ($H_1$ pixels are $\sim$10$\times$ smaller than $H_0$'s), and the CNN has no coordinate channels: an earlier ``CoordConv'' design destabilised training and has been retired.
  \item \textbf{An asymmetry to keep in mind.} For parameter estimation both feature sets get $n(x)$. For classification the PH models do \emph{not}: they were configured as a topology-only baseline, while the classical models always receive $n(x)$. Since the models' parameters are defined in units of the mean point spacing, a persistence diagram without $n(x)$ cannot be converted to those units at all. Part of the classification gap below is that missing column, not the features.
  \item \textbf{Training:} the same for all: early stopping on validation loss, 10 seeds (9371--9380).
\end{itemize}

% ============================================================================
\section{Parameter estimation (inference)}

\begin{itemize}
  \item \textbf{Metric:} normalised loss = mean squared error of the log-parameters, standardised with the training statistics, averaged over the model's parameters. Lower is better; $\approx 1$ means predicting the training mean. Values are \emph{not} comparable across models.
  \item \textbf{Parameters estimated:} Thomas $(\kappa,\mu,\sigma)$; nested $(\kappa,\mu_1,\mu_2,\sigma_1,\sigma)$; Mat\'ern II $(\lambda_p,R)$; LGCP $(\nbar,\sigma^2,s)$.
  \item \textbf{Bold} marks the single best model per family.
\end{itemize}

\begin{longtable}{@{}ll c cccc@{}}
\caption{Parameter estimation: normalised loss (lower is better), mean $\pm$ SD over 10 seeds, with the breakdown by distance from Poisson. All 15 non-empty subsets of $\{L,F,G,J\}$ are shown. \textbf{Bold} marks the single best model per family.}
\label{tab:params}\\
\toprule
Family & Model & Overall & \texttt{below} & \texttt{weak} & \texttt{moderate} & \texttt{strong}\\
\midrule
\endfirsthead
\multicolumn{7}{l}{\small\slshape (Table~\ref{tab:params} continued)}\\
\toprule
Family & Model & Overall & \texttt{below} & \texttt{weak} & \texttt{moderate} & \texttt{strong}\\
\midrule
\endhead
\midrule
\multicolumn{7}{r}{\small\slshape continued on next page}\\
\endfoot
\bottomrule
\endlastfoot
Thomas & $L$ & 0.1690 $\pm$ 0.0082 & 0.282 & 0.234 & 0.115 & 0.045\\
 & $F$ & 0.2584 $\pm$ 0.0151 & 0.355 & 0.375 & 0.228 & 0.075\\
 & $G$ & 0.4735 $\pm$ 0.0191 & 0.367 & 0.504 & 0.593 & 0.429\\
 & $J$ & 0.3287 $\pm$ 0.0185 & 0.351 & 0.439 & 0.376 & 0.149\\
 & $L{+}F$ & 0.1699 $\pm$ 0.0067 & 0.290 & 0.236 & 0.112 & 0.042\\
 & $L{+}G$ & 0.1707 $\pm$ 0.0069 & 0.286 & 0.238 & 0.114 & 0.044\\
 & $L{+}J$ & 0.1698 $\pm$ 0.0078 & 0.286 & 0.235 & 0.113 & 0.045\\
 & $F{+}G$ & 0.2366 $\pm$ 0.0076 & 0.321 & 0.341 & 0.210 & 0.075\\
 & $F{+}J$ & 0.2587 $\pm$ 0.0135 & 0.334 & 0.376 & 0.240 & 0.085\\
 & $G{+}J$ & 0.2611 $\pm$ 0.0104 & 0.332 & 0.363 & 0.245 & 0.104\\
 & $L{+}F{+}G$ & \B{0.1677 $\pm$ 0.0090} & \B{0.284} & \B{0.236} & \B{0.109} & \B{0.042}\\
 & $L{+}F{+}J$ & 0.1730 $\pm$ 0.0108 & 0.285 & 0.243 & 0.118 & 0.047\\
 & $L{+}G{+}J$ & 0.1745 $\pm$ 0.0103 & 0.283 & 0.244 & 0.122 & 0.049\\
 & $F{+}G{+}J$ & 0.2560 $\pm$ 0.0198 & 0.342 & 0.367 & 0.232 & 0.083\\
 & $L{+}F{+}G{+}J$ & 0.1697 $\pm$ 0.0076 & 0.283 & 0.240 & 0.113 & 0.043\\
 & PH image $H_0$ & 0.2487 $\pm$ 0.0082 & 0.328 & 0.357 & 0.228 & 0.081\\
 & PH image $H_1$ & 0.3701 $\pm$ 0.0199 & 0.444 & 0.451 & 0.381 & 0.204\\
 & PH image $H_0{+}H_1$ & 0.2308 $\pm$ 0.0167 & 0.320 & 0.329 & 0.202 & 0.072\\
\midrule
Nested & $L$ & 0.2333 $\pm$ 0.0108 & 0.338 & 0.271 & 0.180 & 0.145\\
 & $F$ & 0.3242 $\pm$ 0.0145 & 0.372 & 0.399 & 0.288 & 0.238\\
 & $G$ & 0.2528 $\pm$ 0.0114 & 0.331 & 0.305 & 0.206 & 0.169\\
 & $J$ & 0.2695 $\pm$ 0.0157 & 0.359 & 0.331 & 0.217 & 0.171\\
 & $L{+}F$ & 0.2361 $\pm$ 0.0164 & 0.334 & 0.282 & 0.184 & 0.144\\
 & $L{+}G$ & 0.2333 $\pm$ 0.0079 & 0.344 & 0.280 & 0.173 & 0.136\\
 & $L{+}J$ & 0.2299 $\pm$ 0.0065 & 0.337 & 0.273 & 0.174 & 0.136\\
 & $F{+}G$ & 0.2486 $\pm$ 0.0149 & 0.343 & 0.299 & 0.190 & 0.162\\
 & $F{+}J$ & 0.2555 $\pm$ 0.0114 & 0.346 & 0.308 & 0.206 & 0.162\\
 & $G{+}J$ & 0.2525 $\pm$ 0.0132 & 0.345 & 0.299 & 0.199 & 0.167\\
 & $L{+}F{+}G$ & \B{0.2287 $\pm$ 0.0130} & \B{0.331} & \B{0.274} & \B{0.175} & \B{0.135}\\
 & $L{+}F{+}J$ & 0.2347 $\pm$ 0.0138 & 0.337 & 0.283 & 0.184 & 0.135\\
 & $L{+}G{+}J$ & 0.2397 $\pm$ 0.0141 & 0.345 & 0.289 & 0.185 & 0.140\\
 & $F{+}G{+}J$ & 0.2501 $\pm$ 0.0147 & 0.336 & 0.302 & 0.204 & 0.159\\
 & $L{+}F{+}G{+}J$ & 0.2333 $\pm$ 0.0108 & 0.338 & 0.271 & 0.180 & 0.145\\
 & PH image $H_0$ & 0.2867 $\pm$ 0.0125 & 0.371 & 0.361 & 0.238 & 0.177\\
 & PH image $H_1$ & 0.4143 $\pm$ 0.0172 & 0.429 & 0.452 & 0.401 & 0.375\\
 & PH image $H_0{+}H_1$ & 0.3022 $\pm$ 0.0126 & 0.392 & 0.379 & 0.250 & 0.188\\
\midrule
Mat\'ern II & $L$ & 0.0137 $\pm$ 0.0005 & 0.027 & 0.013 & 0.008 & 0.007\\
 & $F$ & 0.0911 $\pm$ 0.0057 & 0.105 & 0.132 & 0.107 & 0.020\\
 & $G$ & \B{0.0135 $\pm$ 0.0002} & \B{0.028} & \B{0.013} & \B{0.007} & \B{0.006}\\
 & $J$ & 0.0225 $\pm$ 0.0013 & 0.045 & 0.026 & 0.012 & 0.007\\
 & $L{+}F$ & 0.0136 $\pm$ 0.0003 & 0.028 & 0.012 & 0.008 & 0.007\\
 & $L{+}G$ & 0.0141 $\pm$ 0.0006 & 0.028 & 0.014 & 0.008 & 0.007\\
 & $L{+}J$ & 0.0137 $\pm$ 0.0006 & 0.027 & 0.013 & 0.008 & 0.006\\
 & $F{+}G$ & 0.0166 $\pm$ 0.0008 & 0.034 & 0.018 & 0.008 & 0.006\\
 & $F{+}J$ & 0.0180 $\pm$ 0.0009 & 0.036 & 0.020 & 0.010 & 0.007\\
 & $G{+}J$ & 0.0171 $\pm$ 0.0005 & 0.034 & 0.018 & 0.009 & 0.006\\
 & $L{+}F{+}G$ & 0.0140 $\pm$ 0.0003 & 0.027 & 0.014 & 0.008 & 0.006\\
 & $L{+}F{+}J$ & 0.0141 $\pm$ 0.0004 & 0.027 & 0.015 & 0.008 & 0.006\\
 & $L{+}G{+}J$ & 0.0141 $\pm$ 0.0005 & 0.028 & 0.015 & 0.008 & 0.006\\
 & $F{+}G{+}J$ & 0.0175 $\pm$ 0.0008 & 0.036 & 0.018 & 0.009 & 0.007\\
 & $L{+}F{+}G{+}J$ & 0.0137 $\pm$ 0.0005 & 0.027 & 0.013 & 0.008 & 0.007\\
 & PH image $H_0$ & 0.0361 $\pm$ 0.0017 & 0.058 & 0.049 & 0.027 & 0.011\\
 & PH image $H_1$ & 0.2692 $\pm$ 0.0243 & 0.234 & 0.078 & 0.338 & 0.427\\
 & PH image $H_0{+}H_1$ & 0.0373 $\pm$ 0.0022 & 0.061 & 0.050 & 0.027 & 0.012\\
\midrule
LGCP & $L$ & \B{0.1706 $\pm$ 0.0037} & \B{0.217} & \B{0.174} & \B{0.156} & \B{0.135}\\
 & $F$ & 0.2038 $\pm$ 0.0036 & 0.227 & 0.205 & 0.196 & 0.187\\
 & $G$ & 0.2303 $\pm$ 0.0023 & 0.245 & 0.217 & 0.220 & 0.239\\
 & $J$ & 0.2203 $\pm$ 0.0062 & 0.247 & 0.223 & 0.206 & 0.205\\
 & $L{+}F$ & 0.1710 $\pm$ 0.0032 & 0.220 & 0.175 & 0.154 & 0.135\\
 & $L{+}G$ & 0.1713 $\pm$ 0.0037 & 0.218 & 0.174 & 0.155 & 0.137\\
 & $L{+}J$ & 0.1733 $\pm$ 0.0031 & 0.218 & 0.177 & 0.159 & 0.139\\
 & $F{+}G$ & 0.2070 $\pm$ 0.0041 & 0.238 & 0.212 & 0.200 & 0.178\\
 & $F{+}J$ & 0.2084 $\pm$ 0.0029 & 0.236 & 0.214 & 0.200 & 0.184\\
 & $G{+}J$ & 0.2086 $\pm$ 0.0033 & 0.237 & 0.214 & 0.200 & 0.183\\
 & $L{+}F{+}G$ & 0.1724 $\pm$ 0.0039 & 0.219 & 0.180 & 0.157 & 0.134\\
 & $L{+}F{+}J$ & 0.1738 $\pm$ 0.0034 & 0.222 & 0.181 & 0.157 & 0.135\\
 & $L{+}G{+}J$ & 0.1741 $\pm$ 0.0033 & 0.220 & 0.180 & 0.159 & 0.138\\
 & $F{+}G{+}J$ & 0.2080 $\pm$ 0.0042 & 0.238 & 0.212 & 0.199 & 0.183\\
 & $L{+}F{+}G{+}J$ & \B{0.1706 $\pm$ 0.0037} & \B{0.217} & \B{0.174} & \B{0.156} & \B{0.135}\\
 & PH image $H_0$ & 0.1877 $\pm$ 0.0056 & 0.230 & 0.191 & 0.173 & 0.158\\
 & PH image $H_1$ & 0.2264 $\pm$ 0.0043 & 0.254 & 0.213 & 0.220 & 0.219\\
 & PH image $H_0{+}H_1$ & 0.1834 $\pm$ 0.0038 & 0.222 & 0.192 & 0.171 & 0.149\\
\end{longtable}


\begin{itemize}
  \item \textbf{A classical model wins on every family.} Thomas $L{+}F{+}G$ 0.168, nested $L{+}F{+}G$ 0.229, Mat\'ern~II $G$ 0.014, LGCP $L{+}F{+}G{+}J$ 0.171. As before, several subsets tie within about one seed SD, so read the winners as a tie among the $L$-containing subsets rather than a ranking --- with the exception of Mat\'ern~II, where $G$ alone is genuinely the right summary because a hard core is exactly what $G$ measures.
  \item \textbf{PH trails on every family, and the margin narrows as the signal grows.} On Thomas the best PH model is 0.231 against 0.168; the gap is 0.036 in \texttt{below} but 0.030 in \texttt{strong}, where PH reaches 0.072 against 0.042. On LGCP the two are close throughout (0.183 vs.\ 0.171 overall; 0.149 vs.\ 0.135 in \texttt{strong}).
  \item \textbf{Mat\'ern~II is where PH does worst in relative terms} (0.036 vs.\ 0.014, a factor of 2.6). This is a filtration problem rather than a topology problem: the model is defined by a hard core, and DTM with $k=5$ averages over the five nearest neighbours, which smooths away the minimum-interpoint distance that defines it. An alpha or \v{C}ech filtration puts that distance directly into the $H_0$ death times. It is the clearest single opportunity in the table.
  \item \textbf{$H_1$ is much weaker than $H_0$ everywhere} (Thomas 0.370, nested 0.414, Mat\'ern~II 0.269, LGCP 0.226). On Mat\'ern~II its loss is \emph{not} monotone in $\dtil$ --- it gets worse as the pattern becomes more structured (0.234 in \texttt{below} rising to 0.427 in \texttt{strong}) --- which no other model does and which is worth understanding before $H_1$ is used further.
  \item \textbf{Adding $H_1$ to $H_0$ no longer helps uniformly.} It gains on Thomas (0.231 vs.\ 0.249) and LGCP (0.183 vs.\ 0.188), and now \emph{loses} on nested (0.302 vs.\ 0.287) and Mat\'ern~II (0.037 vs.\ 0.036). On the previous headline set the nested comparison came out the other way; the ranking was inside the noise then and is inside it now.
  \item Loss falls monotonically with $\dtil$ for every model and every family except PH $H_1$ on Mat\'ern~II, which is a useful check that the bands are cutting the axis they claim to.
\end{itemize}

% ============================================================================
\section{Classification}

\begin{itemize}
  \item \textbf{Task:} 5-way classification (Poisson, Thomas, nested, Mat\'ern II, LGCP), all models trained on the same pooled training set.
  \item \textbf{Metric:} accuracy (higher is better; chance is 0.20 in every band, by construction --- see the note on Poisson above).
  \item Recall columns are computed over the whole test set.
\end{itemize}

\begin{table}[ht]
\centering\footnotesize
\begin{tabular}{@{}l c cccc ccccc@{}}
\toprule
 & & \multicolumn{4}{c}{By distance from Poisson} & \multicolumn{5}{c}{Recall per family}\\
\cmidrule(lr){3-6}\cmidrule(l){7-11}
Model & Overall & \texttt{below} & \texttt{weak} & \texttt{mod.} & \texttt{strong} & Poisson & Thomas & Nested & Mat\'ern & LGCP\\
\midrule
$L$ & 0.709 $\pm$ 0.003 & 0.470 & 0.672 & 0.804 & 0.890 & 0.813 & 0.450 & 0.573 & 0.946 & 0.763\\
$F$ & 0.578 $\pm$ 0.004 & 0.326 & 0.533 & 0.667 & 0.785 & 0.795 & 0.333 & 0.404 & 0.709 & 0.648\\
$G$ & 0.587 $\pm$ 0.002 & 0.382 & 0.549 & 0.652 & 0.765 & 0.799 & 0.372 & 0.307 & 0.919 & 0.539\\
$J$ & 0.604 $\pm$ 0.004 & 0.395 & 0.566 & 0.680 & 0.777 & 0.764 & 0.392 & 0.340 & 0.923 & 0.603\\
$L{+}F$ & 0.706 $\pm$ 0.003 & 0.467 & 0.672 & 0.800 & 0.883 & 0.823 & 0.476 & 0.574 & 0.943 & 0.713\\
$L{+}G$ & 0.705 $\pm$ 0.004 & 0.470 & 0.667 & 0.795 & 0.887 & 0.798 & 0.475 & 0.567 & 0.951 & 0.733\\
$L{+}J$ & 0.717 $\pm$ 0.003 & 0.471 & 0.684 & 0.815 & 0.898 & 0.824 & 0.475 & 0.603 & 0.941 & 0.741\\
$F{+}G$ & 0.652 $\pm$ 0.004 & 0.412 & 0.607 & 0.737 & 0.851 & 0.795 & 0.404 & 0.472 & 0.927 & 0.660\\
$F{+}J$ & 0.638 $\pm$ 0.008 & 0.407 & 0.591 & 0.716 & 0.837 & 0.778 & 0.396 & 0.436 & 0.924 & 0.654\\
$G{+}J$ & 0.631 $\pm$ 0.006 & 0.410 & 0.594 & 0.716 & 0.804 & 0.771 & 0.413 & 0.392 & 0.926 & 0.653\\
$L{+}F{+}G$ & \B{0.719 $\pm$ 0.003} & \B{0.473} & \B{0.684} & \B{0.819} & \B{0.899} & 0.813 & 0.478 & 0.613 & 0.945 & 0.747\\
$L{+}F{+}J$ & 0.717 $\pm$ 0.006 & 0.471 & 0.682 & 0.820 & 0.895 & 0.811 & 0.474 & 0.611 & 0.943 & 0.747\\
$L{+}G{+}J$ & 0.715 $\pm$ 0.004 & 0.476 & 0.679 & 0.812 & 0.893 & 0.804 & 0.450 & 0.623 & 0.948 & 0.750\\
$F{+}G{+}J$ & 0.639 $\pm$ 0.004 & 0.407 & 0.598 & 0.716 & 0.834 & 0.793 & 0.424 & 0.444 & 0.915 & 0.617\\
$L{+}F{+}G{+}J$ & 0.716 $\pm$ 0.004 & 0.466 & 0.680 & 0.818 & 0.901 & 0.831 & 0.457 & 0.598 & 0.940 & 0.756\\
PH image $H_0$ & 0.661 $\pm$ 0.004 & 0.356 & 0.639 & 0.776 & 0.872 & 0.825 & 0.371 & 0.572 & 0.807 & 0.729\\
PH image $H_1$ & 0.414 $\pm$ 0.004 & 0.257 & 0.332 & 0.434 & 0.634 & 0.668 & 0.190 & 0.294 & 0.382 & 0.537\\
PH image $H_0{+}H_1$ & 0.666 $\pm$ 0.003 & 0.358 & 0.643 & 0.784 & 0.878 & 0.832 & 0.389 & 0.586 & 0.799 & 0.725\\
\bottomrule
\end{tabular}
\caption{Classification accuracy, mean $\pm$ SD over 10 seeds, on the stratified test set. All 15 non-empty subsets of $\{L,F,G,J\}$; $L$, $L{+}F{+}G{+}J$, $G$ and $F$ report whichever $F/G/J$ radius grid won. \textbf{Bold}: best overall model.}
\label{tab:classify}
\end{table}


\begin{itemize}
  \item \textbf{\bestclassicalname\ is the best classifier overall} (\bestclassical), with the other $L$-containing subsets close behind. The published $L$-only baseline reaches \lonly.
  \item \textbf{The PH deficit is entirely in the undetectable band.} $H_0{+}H_1$ scores \bestph\ overall, but the four bands tell a different story: \phbelow\ / \phweak\ / \phmoderate\ / \phstrong\ against the $L$-only baseline's \lonlyxbelow\ / \lonlyxweak\ / \lonlyxmoderate\ / \lonlyxstrong. PH \emph{beats} $L$ alone in \texttt{moderate} and \texttt{strong} --- by 1.8 and 3.7 points --- and loses 11.6 points in \texttt{below}. Since \texttt{below} is by definition the band where the classical test cannot reject Poisson, this is the one place where an aggregate that mixes the bands is actively misleading about the features.
  \item \textbf{Mat\'ern~II recall is the single largest class gap:} 0.799 for PH against 0.945 for the best classical model. That one class accounts for roughly two thirds of the overall deficit, and the cause is the same DTM smoothing described above.
  \item \textbf{$H_1$ alone (0.414) is well above chance but clearly the weakest classifier}; adding it to $H_0$ is worth about half a point (\bestph\ vs.\ 0.661).
  \item \textbf{Read the PH column as a lower bound.} These runs are single-scale ($k=5$) and, unlike every classical model in the table, do not receive $n(x)$. Both are fixable and neither has been tried yet.
\end{itemize}

% ============================================================================
\section{Next steps}

In rough order of expected value per unit of compute:

\begin{itemize}
  \item \textbf{Give the classification PH models $n(x)$}, and run the $H_0{+}H_1$ configuration properly. Retraining only --- the diagrams already exist. This removes the one clearly unfair asymmetry in Table~\ref{tab:classify}.
  \item \textbf{Run the multi-scale filtration ($k=5,10,15$)} that the configs already specify and that won in the earlier report. Needs the $k=10,15$ diagrams computed first. Nested Thomas is the family with the strongest prior reason to benefit, since its two cluster scales are separated by construction.
  \item \textbf{Add an alpha or \v{C}ech channel for Mat\'ern~II.} DTM-$k$ is structurally blind to a hard core; this is the largest single gap in both tables.
  \item \textbf{Report the bands under $G$ and $F$ as well as $L$.} $\dtil$ is built from $L$ --- the competitor's own statistic --- so conditioning on it pins $L$'s signal-to-noise within a band while leaving a topological method's free to vary. The null tables for $G$ and $F$ already exist, so this is a manifest rebuild and a re-score, with no new training.
  \item \textbf{Then:} Betti curves and persistence landscapes; the non-learned classical references (minimum-contrast fits, envelope tests); and, if the high-$\dtil$ story holds up, patterns generated specifically above $\dtil = 8$, where the current pool is thin and uneven across families.
\end{itemize}

\vspace{1em}
\hrule
\vspace{0.5em}
{\small Tables in this document are generated from \texttt{results/testset\_summary.json} by \texttt{scripts/status\_tables.py}; the summary itself is written by \texttt{scripts/rescore\_results.py}. To refresh: re-run both, then recompile. The test set is defined in \texttt{docs/test\_set.md} and built by \texttt{scripts/build\_testset.py}.}

\end{document}
